\documentclass[11pt,a4paper]{article}

\usepackage[T1]{fontenc}
\usepackage[utf8]{inputenc}
\usepackage{lmodern}
\usepackage{amsmath,amssymb,amsthm}
\usepackage{mathtools}
\usepackage{booktabs}
\usepackage{enumitem}
\usepackage[margin=2.7cm]{geometry}
\usepackage{xcolor}
\usepackage[colorlinks=true,linkcolor=blue!50!black,citecolor=green!40!black,urlcolor=blue!50!black]{hyperref}

\newtheorem{theorem}{Theorem}[section]
\newtheorem{lemma}[theorem]{Lemma}
\newtheorem{proposition}[theorem]{Proposition}
\newtheorem{corollary}[theorem]{Corollary}
\theoremstyle{definition}
\newtheorem{definition}[theorem]{Definition}
\newtheorem{example}[theorem]{Example}
\theoremstyle{remark}
\newtheorem{remark}[theorem]{Remark}

\newcommand{\isa}[1]{\mbox{\small\texttt{#1}}}
\newcommand{\Var}{\mathsf{Var}}
\newcommand{\Fun}{\mathsf{Fun}}
\newcommand{\vars}{\mathcal{V}\mathit{ar}}
\newcommand{\eqAC}{\approx}
\newcommand{\eqE}{=_{E}}
\newcommand{\nf}{\mathit{nf}}
\newcommand{\args}{\mathit{args}}
\newcommand{\mkac}{\mathit{mk}}
\newcommand{\targs}{\mathit{targs}}
\newcommand{\fsub}{\mathit{fsub}}
\newcommand{\Sub}{\mathit{Sub}}
\newcommand{\key}{\mathit{key}}
\newcommand{\VNAC}{V_{\mathrm{NAC}}}
\newcommand{\Vtwo}{V_{>1}}
\newcommand{\AS}{\mathit{AS}}
\newcommand{\psize}{\mathit{size}}
\newcommand{\N}{\mathbb{N}}
\newcommand{\mset}[1]{\{\!\{#1\}\!\}}
\newcommand{\sides}{\mathit{sides}}

\title{A Verified, Terminating and Complete\\ First-Order AC-Unification Algorithm\\ in Isabelle/HOL}
\author{Guilherme Borges Brand\~ao}
\date{}

\emergencystretch=3em
\begin{document}
\maketitle

\begin{abstract}
  We present a formalization in Isabelle/HOL of Stickel's algorithm for unification modulo
  associativity and commutativity (AC), in the form certified in PVS by Ayala-Rinc\'on,
  Fern\'andez, Silva, Kutsia and Nantes-Sobrinho. The signature may contain any number of AC
  symbols together with free symbols. The algorithm is parameterized by a solver for homogeneous
  linear Diophantine systems, specified only by soundness and by the property that its output
  generates all solutions; we instantiate it with the algorithm of Contejean and Devie, whose
  correctness and termination are also formalized. We prove that the algorithm always
  terminates, that every substitution it computes is an AC-unifier, and that every AC-unifier is,
  modulo AC, an instance of a computed one. Termination is established with the lexicographic
  measure of the PVS development; the completeness proof replaces the renaming arguments of the
  paper proofs by a simple invariant on substitutions. The development consists of about 6\,200
  lines of Isabelle and contains no unproved statements.
\end{abstract}

\tableofcontents

%======================================================================
\section{Introduction}
%======================================================================

Unification modulo an equational theory $E$ asks, given terms $s$ and $t$, for substitutions
$\sigma$ with $s\sigma =_E t\sigma$. For the theory AC of an associative and commutative symbol
$f$,
\[
  f(x, f(y, z)) \approx f(f(x, y), z), \qquad f(x, y) \approx f(y, x),
\]
unification is decidable and finitary: every unifiable problem has a finite complete set of
unifiers. The classical algorithm is due to Stickel~\cite{Stickel81}; its termination for
arbitrary (non-ground, nested) problems was first proved by Fages~\cite{Fages87}. The central idea
is that an equation between two sums
\[
  f(s_1, \dots, s_m) \approx^? f(t_1, \dots, t_n)
\]
can be solved by counting: if $\sigma$ is a unifier, each argument $s_i\sigma$ and $t_j\sigma$ is a
sum of ``atoms'', and for each atom the number of times it occurs on both sides must coincide. This
gives a homogeneous linear Diophantine equation, whose minimal solutions describe the most general
way to split arguments into sums of fresh variables.

Ayala-Rinc\'on et al.~\cite{AyalaRincon22,AyalaRincon24} certified a version of Stickel's algorithm
in PVS, including a termination proof based on a lexicographic measure. The Diophantine equations
there are solved by a bounded enumeration of candidate solutions. This document describes an
Isabelle/HOL formalization of the same algorithm with the following features.

\begin{itemize}[leftmargin=*]
  \item The signature contains any number of AC symbols, given by a predicate $\mathit{ac}$, and
        arbitrary free symbols. No well-formedness condition on terms is needed: an AC symbol used
        with a number of arguments other than two simply behaves as a free symbol.
  \item The Diophantine solver is a parameter (an Isabelle locale). We instantiate it with the
        algorithm of Contejean and Devie~\cite{BCD90,CD94}, which computes exactly the minimal
        solutions; its correctness and termination are formalized as part of the development.
  \item We prove termination, soundness and completeness of the algorithm on all inputs.
\end{itemize}

The formalization builds on the AFP entries \emph{First-Order Terms}~\cite{AFP-FOT},
\emph{Homogeneous Linear Diophantine Equations}~\cite{AFP-HLDE} and \emph{Farkas' Lemma}~\cite{AFP-Farkas}.
Section~\ref{sec:prelim} fixes the notation and defines equational theories;
Section~\ref{sec:ac} studies AC equality through normal forms;
Section~\ref{sec:dioph} presents the interface to Diophantine solvers and the Contejean--Devie
algorithm; Section~\ref{sec:alg} describes the unification algorithm;
Sections~\ref{sec:sound}, \ref{sec:complete} and~\ref{sec:term} prove soundness, completeness and
termination; Section~\ref{sec:disc} discusses the differences with the PVS development and the
structure of the formalization.

%======================================================================
\section{Preliminaries}\label{sec:prelim}
%======================================================================

\paragraph{Terms and substitutions.}
Terms are those of the AFP entry \emph{First-Order Terms}: a term is either a variable $\Var\,x$
or an application $\Fun\,f\,[t_1,\dots,t_n]$ of a function symbol to a list of terms. Symbols have
no fixed arity. We write $f(t_1, \dots, t_n)$ and simply $x$ for $\Var\,x$ when no confusion
arises. Variables are natural numbers in the algorithm. A substitution is a function $\sigma$ from
variables to terms; $t\sigma$ is its application to $t$, and $\sigma\circ\tau$ denotes the
composition with $t(\sigma\circ\tau) = (t\sigma)\tau$. We write $\vars(t)$ for the set of
variables of $t$, $\{x \mapsto u\}$ for the substitution that maps $x$ to $u$ and is the identity
elsewhere, and $|t|$ for the size of $t$.

A \emph{unification problem} is a list of pairs $s \approx^? t$, written $P$. We write
$\sides(P)$ for the set of terms occurring as a side of an equation in $P$, $\vars(P)$ for the
set of its variables and $P\sigma$ for the list obtained by applying $\sigma$ to all sides.

\paragraph{Equational theories.}
A set of identities $E$ is a set of pairs of terms. The equational theory $=_E$ is the least
congruence that is closed under substitution and contains $E$. We define it inductively by the
rules
\[
  \frac{(l, r)\in E}{l\sigma =_E r\sigma}
  \qquad
  \frac{}{t =_E t}
  \qquad
  \frac{s =_E t}{t =_E s}
  \qquad
  \frac{s =_E t \quad t =_E u}{s =_E u}
  \qquad
  \frac{s_i =_E t_i \ (1\le i\le n)}{f(s_1,\dots,s_n) =_E f(t_1,\dots,t_n)}
\]
The identities may be stated over a type of variables different from that of the terms; in
practice we give the identities of AC by \emph{all their instances}, so the axiom rule is
only used with the identity substitution.

\begin{lemma}[basic properties of $=_E$]\label{lem:eqth}
  The following hold.
  \begin{enumerate}[nosep]
    \item The relation $=_E$ is an equivalence (\isa{eq\_theory\_equivp}).
    \item If $s =_E t$ then $s\sigma =_E t\sigma$ for every substitution $\sigma$
          (\isa{eq\_theory\_subst}).
    \item If $\sigma(x) =_E \tau(x)$ for all $x \in \vars(t)$, then $t\sigma =_E t\tau$
          (\isa{eq\_theory\_subst\_cong}); combining with 2, if moreover $s =_E t$ then
          $s\sigma =_E t\tau$ (\isa{eq\_theory\_subst\_both}).
    \item If $E \subseteq E'$ then $s =_E t$ implies $s =_{E'} t$ (\isa{eq\_theory\_mono}).
  \end{enumerate}
\end{lemma}

The proof of 2 is by rule induction; in the axiom case, $(l\tau)\sigma = l(\tau\circ\sigma)$
is again an instance of the same identity. Item 4 shows, for instance, that AC equality is
contained in ACh equality, the theory obtained by adding a homomorphism.

A generic locale \isa{equational\_theory} packages these facts, so that each concrete theory
obtains them under a qualified name (for example \isa{ac.subst} for AC).

%======================================================================
\section{AC equality}\label{sec:ac}
%======================================================================

\subsection{Several AC symbols}

The signature is described by a predicate $\mathit{ac}$ on function symbols. For each $f$ with
$\mathit{ac}(f)$ the identities are all instances of associativity and commutativity:
\[
  \mathit{AC}_f = \{ (f(s, f(t, u)),\ f(f(s, t), u)) \} \cup \{ (f(s, t),\ f(t, s)) \},
\]
and AC equality $\eqAC$ is the equational theory of $\bigcup_{\mathit{ac}(f)} \mathit{AC}_f$.
An \emph{AC application} of $f$ is a term $f(a, b)$ with $\mathit{ac}(f)$. The axioms only
mention binary applications, so an AC symbol applied to a list of length different from two is
treated exactly as a free symbol; no restriction on the input terms is needed.

\subsection{Normal forms}

To decide $\eqAC$ and to reason about it, we use a canonical representation in which nested
AC applications are flattened and their arguments are collected in a multiset.

\begin{definition}[normal forms]
  Normal forms are given by the datatype
  \[
    n ::= \mathsf{NV}\,x \mid \mathsf{NF}\,f\,[n_1,\dots,n_k] \mid \mathsf{NA}\,f\,M
  \]
  where $M$ is a finite multiset of normal forms. For an AC symbol $f$, the \emph{$f$-arguments}
  of a normal form are $\nf_f(\mathsf{NA}\,f\,M) = M$ and $\nf_f(n) = \mset{n}$ otherwise. The
  normal form of a term is
  \[
    \nf(x) = \mathsf{NV}\,x, \qquad
    \nf(f(s,t)) = \mathsf{NA}\,f\,(\nf_f(\nf\,s) + \nf_f(\nf\,t)) \text{ if } \mathit{ac}(f),
  \]
  and $\nf(f(t_1,\dots,t_n)) = \mathsf{NF}\,f\,[\nf\,t_1,\dots,\nf\,t_n]$ in all other cases.
\end{definition}

\begin{theorem}[\isa{eq\_AC\_iff\_nf}]\label{thm:nf}
  $s \eqAC t$ if and only if $\nf(s) = \nf(t)$.
\end{theorem}

\begin{proof}[Proof sketch]
  For ``only if'' it suffices to check that both sides of every axiom instance have the same
  normal form, and that $\nf$ is compatible with the congruence rule; the multiset union makes
  associativity and commutativity disappear. For ``if'' we proceed by induction on the size of
  $s$. The interesting case is that of two AC applications of $f$ whose multisets of
  $f$-arguments coincide: the lists of $f$-arguments of $s$ and $t$ are then permutations of each
  other up to argument-wise equal normal forms, and the result follows from the induction
  hypothesis and the fact that a sum does not depend, modulo AC, on the order of its summands.
\end{proof}

\begin{example}\label{ex:nf}
  Let $f$ be AC and $g$ free, and write $s + t$ for $f(s,t)$. The terms
  \[
    t_1 = (x + g(a)) + x \qquad\text{and}\qquad t_2 = x + (x + g(a))
  \]
  both have the normal form $\mathsf{NA}\,f\,\mset{\mathsf{NV}\,x, \mathsf{NV}\,x,
  \mathsf{NF}\,g\,[\mathsf{NF}\,a\,[\,]]}$, hence $t_1 \eqAC t_2$. On the other hand, the
  ternary application $f(x, y, z)$ is not an AC application: its normal form is
  $\mathsf{NF}\,f\,[\mathsf{NV}\,x, \mathsf{NV}\,y, \mathsf{NV}\,z]$, so
  $f(x,y,z) \not\eqAC f(y,x,z)$. Normal forms also decide AC equality below free symbols:
  $g(x + y) \eqAC g(y + x)$, but $g(x + y) \not\eqAC g(x) + g(y)$ (the second equation holds
  in ACh, not in AC).
\end{example}

Two consequences are used throughout. First, $\eqAC$ preserves the size of terms
(\isa{eq\_AC\_size}); in particular AC is \emph{regular} and a variable $x$ cannot be
AC-equal to a term properly containing it, so the occurs check is sound modulo AC. Second, the
syntactic flattening of a term is well behaved.

\begin{definition}[flattening]
  For an AC symbol $f$, $\args_f(t)$ is the list of $f$-arguments of $t$:
  $\args_f(f(s,t)) = \args_f(s) \mathbin{@} \args_f(t)$ and $\args_f(t) = [t]$ if $t$ is not an
  AC application of $f$. Conversely, $\mkac_f[t_1,\dots,t_n] = f(t_1, f(t_2, \dots f(t_{n-1},
  t_n)))$ for $n \ge 1$, with $\mkac_f[t] = t$.
\end{definition}

\begin{lemma}\label{lem:mk}
  For an AC symbol $f$:
  \begin{enumerate}[nosep]
    \item $\mkac_f(\args_f\,t) \eqAC t$ (\isa{mk\_ac\_args});
    \item if the lists $xs$ and $ys$ are permutations of each other, $\mkac_f\,xs \eqAC \mkac_f\,ys$
          (\isa{mk\_ac\_perm});
    \item for non-empty $xs$ and $ys$, $\mkac_f(xs \mathbin{@} ys) \eqAC f(\mkac_f\,xs,
          \mkac_f\,ys)$ (\isa{mk\_ac\_append});
    \item $(\mkac_f\,xs)\sigma = \mkac_f(xs\sigma)$, and $\args_f(\mkac_f\,xs) = xs$ if no
          element of $xs$ is an AC application of $f$;
    \item $\nf_f(\nf\,t)$ is the multiset of the normal forms of $\args_f(t)$, and
          $\nf_f(\nf(t\sigma)) = \sum_{u\in\args_f(t)} \nf_f(\nf(u\sigma))$;
    \item no element of $\args_f(t)$ is an AC application of $f$, and each is strictly smaller
          than $t$ when $t$ is an AC application of $f$.
  \end{enumerate}
\end{lemma}

The left commutativity law $f(a, f(b, c)) \eqAC f(b, f(a, c))$ (\isa{ac\_left\_comm}), derived
from the two axioms, is the basic tool behind item 2. Regularity of AC gives the following form
of the occurs check, used in the completeness proof.

\begin{lemma}[occurs check, \isa{size\_subst\_occurs}, \isa{occurs\_not\_unifiable}]\label{lem:occurs}
  If $x \in \vars(t)$ and $t \neq x$, then $|\delta(x)| < |t\delta|$ for every $\delta$, and
  consequently $\delta(x) \not\eqAC t\delta$.
\end{lemma}

%======================================================================
\section{Homogeneous linear Diophantine systems}\label{sec:dioph}
%======================================================================

\subsection{The solver interface}

A system is given by a number $n$ of unknowns and a list $A$ of integer rows; a vector
$y\in\N^n$ is a solution if $r\cdot y = 0$ for every row $r$ of $A$. The unification algorithm
only needs a list of solutions from which every solution can be obtained as a linear combination
with natural coefficients. This is captured by a locale.

\begin{definition}[Diophantine solver]
  A function $\mathit{dsolve}$ returning, for $n$ and $A$, either $\mathit{None}$ or a list $B$ of
  vectors is a \emph{Diophantine solver} if, whenever $\mathit{dsolve}\,n\,A = \mathit{Some}\,B$,
  \begin{enumerate}[nosep]
    \item every $b \in B$ is a solution, and
    \item every solution $y$ is a combination $y = \sum_i k_i b_i$ of the elements of $B$ with
          $k_i \in \N$.
  \end{enumerate}
\end{definition}

The solver may give up (return $\mathit{None}$); this keeps the interface usable for incomplete
solvers. The function that always returns $\mathit{None}$ is a (useless) solver. Completeness of
the unification algorithm requires a \emph{total} solver, i.e., one that never returns
$\mathit{None}$; soundness and termination hold for any solver.

\begin{lemma}[\isa{minimal\_below}, \isa{Solutions\_gen\_by}]\label{lem:minimal}
  Every non-zero solution lies pointwise above a minimal non-zero solution. Consequently, every
  list containing all minimal non-zero solutions generates all solutions.
\end{lemma}

\begin{proof}
  The first part is by induction on the sum of the components. For the second, if $y$ is a
  non-zero solution and $m \le y$ is minimal, then $y - m$ is again a solution (the system is
  homogeneous) with a smaller sum, so $y = m + (y - m)$ and we conclude by induction.
\end{proof}

Hence any algorithm that computes a superset of the minimal solutions is a solver. Besides the
Contejean--Devie algorithm, the development contains two further instances:
\begin{itemize}[nosep,leftmargin=*]
  \item a solver by \emph{successive elimination} built on the AFP solver for one equation: the
        solutions of the first equation are $y = \sum_i \lambda_i b_i$, and the remaining
        equations become equations in the coefficients $\lambda$ (\isa{dsys\_sol},
        \isa{dsys\_gen});
  \item the \emph{bounded enumeration} of PVS as the solver for one equation
        $a \cdot x = b \cdot y$, extended to systems in the same way: all non-zero vectors whose
        entries are below Stickel's bound $\max(m,n)\cdot\max_{i,j}\mathrm{lcm}(a_i, b_j)$, which
        is correct because it dominates Huet's bounds on minimal solutions.
\end{itemize}

\subsection{The algorithm of Contejean and Devie}

Boudet, Contejean and Devie~\cite{BCD90} generalized Fortenbacher's algorithm for one equation to
systems $Ax = 0$. Starting from the unit vectors $e_1,\dots,e_n$, a breadth-first search builds
candidate vectors. A vector $v$ with $Av \neq 0$ may be incremented in coordinate $j$ only if
\[
  (C_p)\qquad Av \cdot Ae_j < 0,
\]
that is, if the increment moves $Av$ towards $0$; candidates that lie above a solution already
found are discarded, solutions are collected, and the remaining candidates are expanded. The
search proceeds level by level: all vectors of the $k$-th frontier have components summing to
$k$.

\begin{example}\label{ex:cd}
  Consider the equation $2x = y + z$, that is, $A = (2\ {-1}\ {-1})$. The unit vectors have
  images $2$, $-1$ and $-1$.
  \begin{itemize}[nosep,leftmargin=*]
    \item Level 1: no unit vector is a solution. From $e_1$ (image $2$) we may increment the
          coordinates with negative image, giving $(1,1,0)$ and $(1,0,1)$; from $e_2$ and $e_3$
          (image $-1$) only the first coordinate, giving the same two vectors.
    \item Level 2: $(1,1,0)$ and $(1,0,1)$ have image $1$; they are expanded to $(1,2,0)$,
          $(1,1,1)$ and $(1,0,2)$.
    \item Level 3: all three vectors are solutions; the frontier becomes empty.
  \end{itemize}
  The result is $[(1,2,0), (1,1,1), (1,0,2)]$, exactly the minimal non-zero solutions.
  Condition $(C_p)$ pruned, for example, $(2,0,0)$, whose image $4$ is farther from $0$.
\end{example}

We formalize this algorithm and its correctness:

\begin{theorem}[Propositions 1 and 2 of \cite{BCD90}]
  If the search stops, the set of vectors it returns is exactly the set of minimal non-zero
  solutions of $Ax = 0$.
\end{theorem}

The proof of termination in~\cite{BCD90} uses compactness of $[0,1]^n$ and rational
approximations. We use instead König's lemma and Gordan's theorem of the alternative, which is a
consequence of Farkas' lemma~\cite{AFP-Farkas}.

\begin{lemma}[Gordan's alternative, \isa{gordan}]
  For every set $J$ of coordinates, either there is a non-zero solution $y \in \N^n$ of $Ay = 0$
  with $y_i = 0$ for $i \notin J$, or there is a rational vector $z$ with $z \cdot Ae_j \ge 1$ for
  every $j \in J$.
\end{lemma}

\begin{lemma}[König's lemma for the search, \isa{inf\_branch}]
  If no level of the search has an empty frontier, there is an infinite sequence
  $b_0, b_1, \dots$ such that each $b_{k+1}$ is obtained from $b_k$ by an increment allowed by
  $(C_p)$ at level $k$.
\end{lemma}

\begin{theorem}[Proposition 3 of \cite{BCD90}]
  The search always stops. Hence the total function \isa{cd\_solve\_total} computes exactly the
  minimal solutions, and $\lambda n\,A.\ \mathit{Some}(\mathit{cd\_solve\_total}\ n\ A)$ is a
  total Diophantine solver.
\end{theorem}

\begin{proof}[Proof sketch]
  If the search does not terminate, the search tree is infinite and finitely branching, so by
  König's lemma it has an infinite branch $b_0, b_1, \dots$. Condition $(C_p)$ implies
  $\lVert A b_k\rVert^2 \le (k+1)\,C$ for a constant $C$. Let $J$ be the set of coordinates that
  are unbounded along the branch. By Gordan's alternative, either some non-zero solution is
  supported on $J$, and then the branch eventually lies above it and would have been cut; or
  there is a vector $z$ with $z \cdot Ae_j \ge 1$ for all $j \in J$, and then $z \cdot Ab_k$ grows
  linearly in $k$, contradicting the bound on $\lVert Ab_k\rVert$.
\end{proof}

%======================================================================
\section{The algorithm}\label{sec:alg}
%======================================================================

The algorithm works on a problem $P$, the substitution $\sigma$ computed so far, and a natural
number $\mathit{nv}$ such that all variables in use are smaller than $\mathit{nv}$; fresh
variables are $\mathit{nv}, \mathit{nv}+1, \dots$. An equation is an \emph{AC pair} if both sides
are AC applications of the same AC symbol. The algorithm follows the structure of the PVS
version: equations that are not AC pairs are solved as in syntactic unification; once only AC
pairs remain, all of them are handled in one \emph{AC round}.

\subsection{Solving one AC equation}

Let $t = f(\dots)$ and $s = f(\dots)$ be AC applications of the same AC symbol $f$, all of whose
variables are below $\mathit{nv}$.

\begin{enumerate}[leftmargin=*]
  \item \emph{Elimination of common arguments.} Let $L = \mathrm{mset}(\args_f\,t) -
        \mathrm{mset}(\args_f\,s)$ and $R = \mathrm{mset}(\args_f\,s) - \mathrm{mset}(\args_f\,t)$.
        If both are empty, the equation holds modulo AC and it is replaced by nothing.
  \item \emph{Grouping.} Let $u_1,\dots,u_k$ be the distinct elements of $L$ with multiplicities
        $a_1,\dots,a_k$, and $v_1,\dots,v_l$ those of $R$ with multiplicities $b_1,\dots,b_l$.
        Write $w_1,\dots,w_{k+l}$ for $u_1,\dots,u_k,v_1,\dots,v_l$.
  \item \emph{Diophantine equation.} Solve
        \[ a_1x_1 + \dots + a_kx_k = b_1y_1 + \dots + b_ly_l \]
        with the solver, and let $B$ be the list of non-zero solutions returned (empty if the
        solver gives up).
  \item \emph{Choice of solutions.} For every sublist $D = [d_0,\dots,d_{p-1}]$ of $B$ that
        satisfies the constraints of Fages: for each $j$, the column sum
        $c_j = \sum_i d_i(j)$ is positive, and $c_j = 1$ if $w_j$ is not a variable, produce
        the equations
        \[
          w_j \approx^? \mkac_f\bigl(\underbrace{z_0,\dots,z_0}_{d_0(j)},\ \dots,\
          \underbrace{z_{p-1},\dots,z_{p-1}}_{d_{p-1}(j)}\bigr)
          \qquad (1 \le j \le k+l),
        \]
        where $z_i = \mathit{nv} + i$ are fresh, and the new bound $\mathit{nv} + p$.
\end{enumerate}

We call this function $\mathit{solve\_ac}$; it returns a list of alternatives $(C, \mathit{nv}')$.

\begin{lemma}[shape of the alternatives, \isa{solve\_ac\_shape}]\label{lem:shape}
  For every alternative $(C, \mathit{nv}')$ of $\mathit{solve\_ac}\ f\ t\ s\ \mathit{nv}$:
  $\mathit{nv} \le \mathit{nv}'$; the left-hand sides of $C$ are distinct elements of
  $\args_f(t) \mathbin{@} \args_f(s)$, hence none is an AC application of $f$; every right-hand
  side is either a fresh variable $z$ with $\mathit{nv} \le z < \mathit{nv}'$ or an AC application
  of $f$ whose $f$-arguments are such variables; and a non-variable left-hand side has a variable
  right-hand side.
\end{lemma}

\begin{example}\label{ex:xxyz}
  Let $+$ be AC and consider $x + x \approx^? y + z$. There are no common arguments;
  $L = \mset{x, x}$ and $R = \mset{y, z}$ give $w = (x, y, z)$ and the equation
  $2X = Y + Z$, whose minimal solutions were computed in Example~\ref{ex:cd}:
  \[
    b_1 = (1,2,0), \qquad b_2 = (1,1,1), \qquad b_3 = (1,0,2).
  \]
  All $w_j$ are variables, so a sublist $D$ is admissible if every column has a positive sum: the
  column of $Y$ needs $b_1$ or $b_2$, and that of $Z$ needs $b_2$ or $b_3$. This leaves five
  choices; with fresh variables $z_0, z_1, z_2$ for the elements of $D$ they give:
  \begin{center}
  \footnotesize
  \begin{tabular}{@{}lll@{}}
    \toprule
    $D$ & equations $C$ & unifier, on $x, y, z$\\
    \midrule
    $[b_2]$ & $x \approx^? z_0,\ y \approx^? z_0,\ z \approx^? z_0$ & $x, y \mapsto z$\\
    $[b_1, b_2]$ & $x \approx^? z_0 + z_1,\ y \approx^? z_0 + z_0 + z_1,\ z \approx^? z_1$ &
       $x \mapsto z_0 + z,\ y \mapsto z_0 + z_0 + z$\\
    $[b_1, b_3]$ & $x \approx^? z_0 + z_1,\ y \approx^? z_0 + z_0,\ z \approx^? z_1 + z_1$ &
       $x \mapsto z_0 + z_1,\ y \mapsto z_0 + z_0,\ z \mapsto z_1 + z_1$\\
    $[b_2, b_3]$ & $x \approx^? z_0 + z_1,\ y \approx^? z_0,\ z \approx^? z_0 + z_1 + z_1$ &
       $x \mapsto y + z_1,\ z \mapsto y + z_1 + z_1$\\
    $[b_1, b_2, b_3]$ & $x \approx^? z_0 + z_1 + z_2,\ y \approx^? z_0 + z_0 + z_1,$ &
       $x \mapsto z_0 + z_1 + z_2,\ y \mapsto z_0 + z_0 + z_1,$\\
       & $z \approx^? z_1 + z_2 + z_2$ & $z \mapsto z_1 + z_2 + z_2$\\
    \bottomrule
  \end{tabular}
  \end{center}
  The last column is the result of $\mathit{inst}$ described below, which binds a fresh variable
  on the right to the old variable on the left (so, in the first row, $z_0 := x$, then
  $x := y$, then $y := z$). Each unifier can be checked by counting: in the third row,
  $x + x$ and $y + z$ both become $z_0 + z_0 + z_1 + z_1$ modulo AC.
\end{example}

\begin{example}[the constraints of Fages]\label{ex:fages}
  Let $+$ be AC, $g$ free and $a$ a constant, and consider $x + g(a) \approx^? y + y$. Now
  $w = (x, g(a), y)$ with the equation $X_1 + X_2 = 2Y$. The search of Contejean and Devie returns
  $b_1 = (2,0,1)$, $b_2 = (1,1,1)$ and $b_3 = (0,2,1)$. Since $g(a)$ is not a variable, its
  column must sum to exactly $1$, which rules out $[b_1]$, $[b_3]$ and every $D$ containing $b_3$.
  The two admissible choices are
  \begin{align*}
    [b_2]:\ &\ x \approx^? z_0,\ g(a) \approx^? z_0,\ y \approx^? z_0,\\
            &\ \text{giving } x \mapsto g(a),\ y \mapsto g(a);\\
    [b_1, b_2]:\ &\ x \approx^? z_0 + z_0 + z_1,\ g(a) \approx^? z_1,\ y \approx^? z_0 + z_1,\\
            &\ \text{giving } x \mapsto z_0 + z_0 + g(a),\ y \mapsto z_0 + g(a).
  \end{align*}
  Without the constraint, $[b_3]$ would produce $g(a) \approx^? z_0 + z_0$, which has no
  solution, since $g(a)$ is not a sum.
\end{example}

\subsection{Instantiation}

The equations produced by $\mathit{solve\_ac}$ are turned into a substitution by
$\mathit{inst}$ (the PVS \emph{instantiateStep}), which processes them from left to right.
An equation $t \approx^? t$ is dropped. An equation $t \approx^? z$ with a variable right-hand
side binds $z := t$, unless $z \in \vars(t)$, in which case the alternative fails; the
right-hand side is tried first because after an AC step it contains the fresh variables. An
equation $x \approx^? s$ with $s$ not a variable binds $x := s$ under the same occurs check. An
equation between two non-variable terms is kept. The result is either a failure or a pair
$(R, \theta)$ of the kept equations, instantiated, and the computed substitution $\theta$.

\begin{lemma}[properties of instantiation]\label{lem:inst}
  Let $\mathit{inst}\ E = (R, \theta)$.
  \begin{enumerate}[nosep]
    \item $\vars(R) \subseteq \vars(E)$, $\vars(\theta(x)) \subseteq \vars(E) \cup \{x\}$, and
          $\theta(x) = x$ for $x \notin \vars(E)$ (\isa{inst\_vars}).
    \item $\theta$ is idempotent: if $w \in \vars(\theta(y))$ then $\theta(w) = w$
          (\isa{inst\_idem}).
    \item For every $y$, either $\theta(y)$ is a variable $w$ with $\theta(w) = w$, or
          $\theta(y) = a\theta$ for some non-variable side $a$ of $E$ (\isa{inst\_img}).
    \item Every side of $R$ is $a\theta$ for some non-variable side $a$ of $E$ (\isa{inst\_res});
          and if $\theta(x) \neq x$ then $x$ is a side of $E$ (\isa{inst\_dom}).
  \end{enumerate}
\end{lemma}

\subsection{One AC round}

$\mathit{ac\_round}\ P\ R\ \sigma\ \mathit{nv}$ processes the AC pairs of $P$ one by one,
accumulating the kept equations in $R$:
\begin{align*}
  \mathit{ac\_round}\ []\ R\ \sigma\ \mathit{nv} &= [(R, \sigma, \mathit{nv})],\\
  \mathit{ac\_round}\ ((t \approx^? s) \# P)\ R\ \sigma\ \mathit{nv}
    &= \bigl[\ r \mid (C, \mathit{nv}') \in \mathit{solve\_ac}\ f\ t\ s\ \mathit{nv},\
       \mathit{inst}\ C = (R', \theta),\\
    &\qquad\quad r \in \mathit{ac\_round}\ (P\theta)\ (R\theta \mathbin{@} R')\
       (\sigma\circ\theta)\ \mathit{nv}' \ \bigr].
\end{align*}

\subsection{The main loop}

\begin{tabbing}
  xx\=xx\=xx\=\kill
  $\mathit{acunif}\ P\ \sigma\ \mathit{nv}$:\\
  \> if $P = []$ then return $[\sigma]$\\
  \> else if some equation of $P$ is not an AC pair, let $t \approx^? s$ be the first one and
     $P_1 = P \setminus \{t \approx^? s\}$:\\
  \>\> \emph{delete:} if $t = s$, return $\mathit{acunif}\ P_1\ \sigma\ \mathit{nv}$;\\
  \>\> \emph{eliminate:} if $t = x$ and $x \notin \vars(s)$, return
       $\mathit{acunif}\ (P_1\{x\mapsto s\})\ (\sigma\circ\{x\mapsto s\})\ \mathit{nv}$;\\
  \>\> \emph{orient and eliminate:} symmetrically if $s = x$, $t$ is not a variable and
       $x \notin\vars(t)$;\\
  \>\> \emph{decompose:} if $t = g(t_1..t_n)$, $s = g(s_1..s_n)$, return
       $\mathit{acunif}\ ([t_i\approx^?s_i]_i \mathbin{@} P_1)\ \sigma\ \mathit{nv}$;\\
  \>\> \emph{fail} (occurs check or clash): return $[]$;\\
  \> else return the concatenation of $\mathit{acunif}\ P'\ \sigma'\ \mathit{nv}'$ for all
     $(P', \sigma', \mathit{nv}') \in \mathit{ac\_round}\ P\ []\ \sigma\ \mathit{nv}$.
\end{tabbing}

Decomposition is also used for AC symbols applied to a number of arguments other than two, which
are not AC pairs.

\begin{example}[a complete run]\label{ex:run}
  Let $+$ be AC, $g$ free and unary, and $a, b$ constants. On $g(x + y) \approx^? g(a + b)$:
  \begin{enumerate}[nosep,leftmargin=*]
    \item The equation is not an AC pair, so it is decomposed into $x + y \approx^? a + b$.
    \item Only AC pairs remain, so an AC round starts. The arguments give $w = (x, y, a, b)$ and
          the equation $X_1 + X_2 = Y_1 + Y_2$, whose minimal solutions are the four vectors
          $p = (1,0,1,0)$, $q = (1,0,0,1)$, $r = (0,1,1,0)$, $s = (0,1,0,1)$, pairing one
          argument on each side. The columns of the constants $a$ and $b$ must sum to $1$, and
          those of $x$ and $y$ must be positive; only $[p, s]$ and $[q, r]$ qualify.
    \item For $[p, s]$ the equations are $x \approx^? z_0$, $y \approx^? z_1$, $a \approx^? z_0$,
          $b \approx^? z_1$; instantiation gives $z_0 := x$, $z_1 := y$, then $x := a$, $y := b$,
          and no equation is kept. Likewise $[q, r]$ gives $x := b$, $y := a$.
    \item Both branches end with the empty problem.
  \end{enumerate}
  Hence $\mathit{ac\_unifiers}$ returns two unifiers, $\{x \mapsto a, y \mapsto b\}$ and
  $\{x \mapsto b, y \mapsto a\}$ (restricted to $x, y$).
\end{example} The set of AC-unifiers of $s$ and $t$ is
\[
  \mathit{ac\_unifiers}\ s\ t = \mathit{acunif}\ [s \approx^? t]\ \mathit{id}\ (1 + \max(\{0\}\cup\vars(s)\cup\vars(t))).
\]

Since termination is not obvious, $\mathit{acunif}$ is defined in Isabelle as a partial function
(\isa{function (domintros)}) together with its domain predicate. Soundness and completeness are
first proved for arguments in the domain, by partial induction; the termination proof of
Section~\ref{sec:term} shows that every argument is in the domain.

%======================================================================
\section{Soundness}\label{sec:sound}
%======================================================================

A substitution $\delta$ \emph{unifies} $P$ if $s\delta \eqAC t\delta$ for every equation
$s \approx^? t$ of $P$.

\begin{lemma}[soundness of the steps]
  The following hold.
  \begin{enumerate}[nosep]
    \item If $\mathit{inst}\ C = (R, \theta)$ and $\delta$ unifies $R$, then $\theta\circ\delta$
          unifies $C$.
    \item If $(C, \mathit{nv}')$ is an alternative of $\mathit{solve\_ac}\ f\ t\ s\ \mathit{nv}$
          and $\delta$ unifies $C$, then $t\delta \eqAC s\delta$.
    \item If $(R', \sigma', \mathit{nv}')$ is a result of
          $\mathit{ac\_round}\ P\ R\ \sigma\ \mathit{nv}$, then $\sigma' = \sigma\circ\theta$ for
          some $\theta$ such that $\theta\circ\tau$ unifies $P \mathbin{@} R$ whenever $\tau$
          unifies $R'$.
  \end{enumerate}
\end{lemma}

Item 2 is proved with normal forms. Removing the common arguments, it suffices to show that
every normal form $e$ occurs equally often as an $f$-argument of the instances of the arguments in
$L$ and in $R$. Each $w_j\delta$ has the same $f$-arguments as the instance of its right-hand
side, a sum of fresh variables $z_i$ taken $d_i(j)$ times. So both counts are sums, over $i$, of
the occurrences of $e$ in $z_i\delta$ weighted by $\sum_{j} a_j d_i(j)$ and by
$\sum_{j} b_j d_i(j)$ respectively, and these weights coincide because each $d_i$ is a
solution of the Diophantine equation. Theorem~\ref{thm:nf} then gives AC equality.

The other steps are handled by simple facts about unification problems:
\begin{lemma}
  The following hold.
  \begin{enumerate}[nosep]
    \item $\tau$ unifies $P\sigma$ if and only if $\sigma\circ\tau$ unifies $P$
          (\isa{unifies\_subst\_eqs});
    \item if $\tau$ unifies $[t_i \approx^? s_i]_{i}$ with lists of equal length, then $\tau$
          unifies $g(t_1..t_n) \approx^? g(s_1..s_n)$ (\isa{unifies\_zip});
    \item if $x \notin \vars(t)$ then $t\{x \mapsto u\} = t$ (\isa{subst\_upd\_not\_in}).
  \end{enumerate}
\end{lemma}

\begin{theorem}[\isa{acunif\_sound}]
  If the computation of $\mathit{acunif}\ P\ \sigma\ \mathit{nv}$ terminates and
  $\theta$ is one of its results, then $\theta = \sigma\circ\tau$ for some $\tau$ that unifies
  $P$.
\end{theorem}

\begin{theorem}[soundness, \isa{ac\_unifiers\_sound}]
  If $\theta \in \mathit{ac\_unifiers}\ s\ t$ then $s\theta \eqAC t\theta$.
\end{theorem}

%======================================================================
\section{Completeness}\label{sec:complete}
%======================================================================

The paper proofs of completeness of Stickel's algorithm reason about renamings of variables and
about substitutions \emph{restricted} to the variables of the input. We use a simpler invariant.

\begin{definition}
  A substitution $\delta$ \emph{AC-fixes} $\theta$ if $\delta(x) \eqAC (\theta(x))\delta$ for every
  variable $x$.
\end{definition}

\begin{lemma}[AC-fixing]\label{lem:fixes}
  The following hold.
  \begin{enumerate}[nosep]
    \item Every $\delta$ AC-fixes the identity substitution (\isa{fixes\_ac\_Var}).
    \item If $\delta(x) \eqAC t\delta$, then $\delta$ AC-fixes $\{x \mapsto t\}$
          (\isa{fixes\_ac\_upd}).
    \item If $\delta$ AC-fixes $\sigma$ and $\theta$, it AC-fixes $\sigma\circ\theta$
          (\isa{fixes\_ac\_compose}).
    \item If $\delta$ AC-fixes $\theta$, then $u\theta\delta \eqAC u\delta$ for every term $u$
          (\isa{fixes\_ac\_term}).
    \item If $\delta$ AC-fixes $\theta$ and unifies $P$, then it unifies $P\theta$
          (\isa{unifies\_subst\_eqs\_fixes}).
  \end{enumerate}
\end{lemma}

Items 2 and 3 cover the variable elimination steps, and item 5 shows that a unifier of the
problem remains a unifier after instantiation. At the end of a successful run, if $\delta$
AC-fixes the computed $\theta$, then $\delta(x) \eqAC \theta(x)\delta$ for all $x$, so $\delta$
is an AC-instance of $\theta$, with $\delta$ itself as the instantiating substitution.

The only steps that create information not present in $\delta$ are the AC steps, which introduce
fresh variables. The invariant is therefore stated for an extension $\delta'$ of $\delta$ that
agrees with $\delta$ below the current bound $\mathit{nv}$. We also maintain the invariant
$\mathit{inv\_vars}$: the variables of the problem are below $\mathit{nv}$, the substitution is
the identity above $\mathit{nv}$, and it maps variables below $\mathit{nv}$ to terms with
variables below $\mathit{nv}$.

\begin{lemma}[completeness of instantiation]\label{lem:inst-complete}
  If $\delta$ unifies $C$, then $\mathit{inst}\ C$ succeeds with some $(R, \theta)$ such that
  $\delta$ unifies $R$ and AC-fixes $\theta$.
\end{lemma}

The occurs check never fails here, by Lemma~\ref{lem:occurs}. For decomposition we use:

\begin{lemma}[\isa{decompose\_unifiable}]
  If $g(t_1..t_n) \approx^? h(s_1..s_m)$ is not an AC pair and is unified by $\delta$, then
  $g = h$, $n = m$, and $\delta$ unifies $t_i \approx^? s_i$ for every $i$.
\end{lemma}

The AC step relies on two facts about the Diophantine equation.

\begin{lemma}[\isa{mult\_vector\_sol}, \isa{basis\_gen}]
  Let $t\delta \eqAC s\delta$ for AC applications $t, s$ of $f$, with $w$, $a$ and $b$ as in the
  AC step. For every normal form $e$, the vector whose $j$-th component is the number of
  occurrences of $e$ among the $f$-arguments of $w_j\delta$ solves $a \cdot X = b \cdot Y$. If
  the solver does not give up, every solution is a combination with natural coefficients of the
  computed basis.
\end{lemma}

\begin{lemma}[completeness of an AC step]\label{lem:solve-complete}
  Assume the solver is total. Let $t$ and $s$ be AC applications of $f$ with variables below
  $\mathit{nv}$, and $t\delta \eqAC s\delta$. Then $\mathit{solve\_ac}\ f\ t\ s\ \mathit{nv}$ has an
  alternative $(C, \mathit{nv}')$ with $\mathit{nv} \le \mathit{nv}'$, variables below
  $\mathit{nv}'$, and a substitution $\delta'$ that agrees with $\delta$ below $\mathit{nv}$ and
  unifies $C$.
\end{lemma}

\begin{proof}[Proof sketch]
  Consider the set $A$ of normal forms $e$ that occur as $f$-arguments of some $w_j\delta$. For
  each $e \in A$, the vector whose $j$-th component is the number of occurrences of $e$ among the
  $f$-arguments of $w_j\delta$ is a solution of the Diophantine equation: this is
  Theorem~\ref{thm:nf} applied to $t\delta \eqAC s\delta$, after removing the common arguments.
  Since the solver is total, it can be written as $\sum_i c_e(i)\, b_i$ with $c_e(i)\in\N$, where
  $b_i$ ranges over the computed basis. Let $D$ be the sublist of basis vectors $b_i$ with
  $c_e(i) > 0$ for some $e\in A$. Choosing for each $e \in A$ a term $\mathit{rep}(e)$ with normal
  form $e$ that is not an AC application of $f$, define the fresh variable $z_d$ of the $d$-th
  element of $D$ by
  \[
    \delta'(z_d) = \mkac_f\bigl(\text{$\mathit{rep}(e)$ repeated $c_e(i_d)$ times, for each } e \in A\bigr),
  \]
  and $\delta' = \delta$ below $\mathit{nv}$. Counting occurrences of each $e$ on both sides shows
  that $\delta'$ unifies every equation $w_j \approx^? \mkac_f(\dots)$, so by
  Theorem~\ref{thm:nf} it unifies $C$. The constraints of Fages are satisfied: each $w_j\delta$ has
  at least one $f$-argument, and a non-variable $w_j$, not being an AC application of $f$, has
  exactly one.
\end{proof}

\begin{example}\label{ex:complete}
  Take $x + x \approx^? y + z$ of Example~\ref{ex:xxyz} and the unifier
  $\delta = \{x \mapsto a + b,\ y \mapsto a + a,\ z \mapsto b + b\}$, with constants $a$, $b$.
  The atoms are $a$ and $b$. Counting them in $(x\delta, y\delta, z\delta)$ gives $(1,2,0) = b_1$
  for $a$ and $(1,0,2) = b_3$ for $b$, so $c_a = (1,0,0)$, $c_b = (0,0,1)$ and $D = [b_1, b_3]$.
  The construction sets $\delta'(z_0) = a$ and $\delta'(z_1) = b$, and $\delta'$ unifies the
  equations $x \approx^? z_0 + z_1$, $y \approx^? z_0 + z_0$, $z \approx^? z_1 + z_1$ of that
  alternative. Instantiation produces $\theta = \{x \mapsto z_0 + z_1,\ y \mapsto z_0 + z_0,\
  z \mapsto z_1 + z_1\}$, and indeed $\delta(x) \eqAC \theta(x)\delta'$ for all variables: e.g.\
  $\theta(x)\delta' = a + b = \delta(x)$. So $\delta$ is an instance of the computed unifier, with
  $\rho = \delta'$.
\end{example}

\begin{lemma}[completeness of a round, \isa{ac\_round\_complete}]
  Assume the solver is total. Let $P$ consist of AC pairs, let $\delta$ unify $P \mathbin{@} R$
  and AC-fix $\sigma$, and let the variable invariant hold. Then some result
  $(R', \sigma', \mathit{nv}')$ of $\mathit{ac\_round}\ P\ R\ \sigma\ \mathit{nv}$ and some
  $\delta'$ agreeing with $\delta$ below $\mathit{nv}$ are such that $\delta'$ unifies $R'$,
  AC-fixes $\sigma'$, and the invariant holds for $R'$, $\sigma'$ and $\mathit{nv}'$.
\end{lemma}

\begin{theorem}[completeness on the domain, \isa{acunif\_complete}]
  Assume the solver is total, the computation of $\mathit{acunif}\ P\ \sigma\ \mathit{nv}$
  terminates, $\delta$ unifies $P$ and AC-fixes $\sigma$, and the variable invariant holds. Then
  some result $\theta$ and some $\delta'$ agreeing with $\delta$ below $\mathit{nv}$ are such
  that $\delta'$ AC-fixes $\theta$.
\end{theorem}

Combining these with the termination theorem of the next section:

\begin{theorem}[completeness, \isa{ac\_unifiers\_complete}]
  Assume the Diophantine solver is total. If $s\delta \eqAC t\delta$, then there are
  $\theta \in \mathit{ac\_unifiers}\ s\ t$ and $\rho$ such that $\delta(x) \eqAC \theta(x)\rho$ for
  every $x \in \vars(s) \cup \vars(t)$.
\end{theorem}

%======================================================================
\section{Termination}\label{sec:term}
%======================================================================

Termination of Stickel's algorithm for nested problems with several AC symbols is the delicate
part. We follow the lexicographic measure of the appendix of the PVS
development~\cite{AyalaRincon24}, adapted to flattened terms.

\subsection{Flattened subterms and the measure}

For a non-variable term $t$, the \emph{flattened arguments} $\targs(t)$ are $\args_f(t)$ if $t$ is
an AC application of $f$, and the usual arguments otherwise. The \emph{key} of $t = g(\dots)$ is
the pair of $g$ and a flag telling whether $t$ is an AC application. The flattened subterms
$\fsub(t)$ are $t$ together with the flattened subterms of its flattened arguments; for a problem
$P$, $\Sub(P)$ is the union of the $\fsub(u)$ for $u \in \sides(P)$.

\begin{definition}[components of the measure]
  For a problem $P$:
  \begin{itemize}[nosep]
    \item $\VNAC(P)$ is the set of variables $x$ such that $x$ is a side of $P$, or $x$ is a
          flattened argument of a term of $\Sub(P)$ that is not an AC application;
    \item $\Vtwo(P)$ is the set of variables that are flattened arguments of two terms of
          $\Sub(P)$ with different keys;
    \item $\AS(P)$, the \emph{admissible subterms}, is the set of non-variable terms $u$ that are
          proper flattened subterms of a side of $P$;
    \item $\psize(P)$ is the sum of the sizes of all sides of $P$.
  \end{itemize}
  The measure of $P$ is the quadruple
  $\bigl(|\VNAC(P)|, |\Vtwo(P)|, |\AS(P)|, \psize(P)\bigr)$, ordered lexicographically.
\end{definition}

All four components are natural numbers, so the lexicographic order is well-founded.

\begin{example}\label{ex:measure}
  Let $+$ be AC and $g$ free, and $P = [\,x + g(y) \approx^? y + z\,]$. The flattened subterms
  are $x + g(y)$, $x$, $g(y)$, $y$, $y + z$ and $z$. No side is a variable; the only term of
  $\Sub(P)$ that is not an AC application and has a variable argument is $g(y)$, so
  $\VNAC(P) = \{y\}$. The variable $y$ is an argument both of $g(y)$, with key $(g,
  \mathit{false})$, and of $y + z$, with key $(+, \mathit{true})$, so $\Vtwo(P) = \{y\}$; $x$ and
  $z$ only occur below $+$. Finally $\AS(P) = \{g(y)\}$. The measure is $(1, 1, 1, \psize(P))$.

  The AC step gives $w = (x, g(y), y, z)$, the equation $X_1 + X_2 = Y_1 + Y_2$, and the basis
  $p = (1,0,1,0)$, $q = (1,0,0,1)$, $r = (0,1,1,0)$, $s = (0,1,0,1)$. The column of $g(y)$ must
  sum to $1$, so exactly one of $r, s$ is chosen, and four alternatives remain:
  \begin{itemize}[nosep,leftmargin=*]
    \item $[q, r]$ and $[p, q, r]$ pair $g(y)$ with $y$ and contain $y \approx^? z_1$ or
          $y \approx^? z_0 + z_2$ with $z_1$, $z_2$ bound to $g(y)$; instantiation fails on the
          occurs check;
    \item $[p, s]$: $x \approx^? z_0$, $g(y) \approx^? z_1$, $y \approx^? z_0$, $z \approx^? z_1$
          gives $\theta = \{x \mapsto y,\ z \mapsto g(y)\}$ on the original variables;
    \item $[p, q, s]$: $x \approx^? z_0 + z_1$, $g(y) \approx^? z_2$, $y \approx^? z_0$,
          $z \approx^? z_1 + z_2$ gives $\theta = \{x \mapsto y + z_1,\ z \mapsto z_1 + g(y)\}$.
  \end{itemize}
  In both successful branches nothing is kept and the round ends with the empty problem, of
  measure $(0,0,0,0)$. Note that $\theta(y) = y$ in both, in accordance with the hypothesis
  of Lemma~\ref{lem:as-step}: the variable $y \in \Vtwo(P)$ is not bound to a non-variable
  term. So $\mathit{ac\_unifiers}$ returns two unifiers.
\end{example}

\subsection{The syntactic steps}

\begin{lemma}[\isa{mless\_delete}, \isa{mless\_elim}, \isa{mless\_decomp}]
  Each step of the main loop on a non-AC pair decreases the measure:
  \begin{enumerate}[nosep]
    \item deletion does not increase $\VNAC$, $\Vtwo$ and $\AS$, and decreases $\psize$;
    \item elimination of $x$ decreases $\VNAC$: the new problem has $\VNAC \subseteq \VNAC(P)
          \setminus\{x\}$, while $x \in \VNAC(P)$ since $x$ is a side;
    \item decomposition of $g(t_1..t_n) \approx^? g(s_1..s_n)$, which is not an AC pair, does
          not increase $\VNAC$, $\Vtwo$ and $\AS$ (the new sides are flattened arguments of a
          non-AC application), and decreases $\psize$.
  \end{enumerate}
  None of these steps introduces new variables.
\end{lemma}

The elimination case uses a characterization of the flattened subterms of an instance: every
$T \in \fsub(v\sigma)$ is either $w\sigma$ for some non-variable $w \in \fsub(v)$, or lies in
$\fsub(\sigma(x))$ for some $x \in \vars(v)$; in the latter case, when $T = \sigma(x)$, either
$v = x$ or $x$ is a flattened argument of some $w \in \fsub(v)$ that is not an AC application of
the same symbol as $\sigma(x)$ (\isa{fsub\_subst}).

\subsection{One AC step}

Consider one step of an AC round: an AC pair $t \approx^? s$ of $f$ is solved by an alternative
$C$ of $\mathit{solve\_ac}$, $\mathit{inst}\ C = (R, \theta)$, and the problem
$S_0 = (t\approx^?s) \# Q$ becomes $S_1 = Q\theta \mathbin{@} R$, where all sides of $Q$ are
non-variable terms with variables below $\mathit{nv}$. The substitution $\theta$ is idempotent,
and its non-trivial bindings have a precise shape: a binding $\theta(x)$ that is not a variable
is either $u\theta$ for a non-variable argument $u$ of $t$ or $s$, or an AC application of $f$
whose arguments are fresh variables fixed by $\theta$ or terms $u\theta$ as before.

\begin{lemma}[\isa{dom\_orig}, \isa{fixed\_low}, \isa{img\_shape}]\label{lem:theta}
  The following hold.
  \begin{enumerate}[nosep]
    \item If $x < \mathit{nv}$ and $\theta(x) \neq x$, then $x$ is a flattened argument of $t$
          or of $s$; in particular $\theta(x) = x$ if $x < \mathit{nv}$ does not occur in $t$ or $s$.
    \item If $\theta(y)$ is not a variable, then either $\theta(y) = u\theta$ for a non-variable
          argument $u$ of $t$ or $s$, or $\theta(y)$ is an AC application of $f$ each of whose
          $f$-arguments is a variable fixed by $\theta$ or such a $u\theta$.
    \item All variables of $S_1$ are fixed by $\theta$, and
          $\vars(S_1) \subseteq \vars(S_0) \cup \{\mathit{nv}, \dots, \mathit{nv}'-1\}$
          (\isa{S1\_fixed}, \isa{S1\_vars}).
  \end{enumerate}
\end{lemma}

Let $\psi(x) = y$ if $\theta(x)$ is the variable $y$, and $\psi(x) = x$ otherwise.

\begin{lemma}[PVS Lemmas 17 and 18]
  $\VNAC(S_1) \subseteq \psi(\VNAC(S_0))$ and $\Vtwo(S_1) \subseteq \psi(\Vtwo(S_0))$.
\end{lemma}

\begin{lemma}[PVS Lemmas 22 and 23]\label{lem:as-step}
  If $\theta(x)$ is a variable for every $x \in \Vtwo(S_0)$, then
  \[
    \AS(S_1) \subseteq \AS(S_0)\theta \qquad\text{and}\qquad \sides(R) \subseteq \AS(S_0)\theta.
  \]
\end{lemma}

\begin{proof}[Proof sketch]
  By induction on the size of $v\theta$, every non-variable proper flattened subterm $T$ of
  $v\theta$, for a non-variable $v$ that is a side or an admissible subterm of $S_0$, belongs to
  $\AS(S_0)\theta$. By the characterization of $\fsub(v\theta)$, either $T = w\theta$ with
  $w$ a proper flattened subterm of $v$, or $T$ lies inside a binding $\theta(x)$. In the second
  case the shape of the bindings reduces the claim to smaller arguments of $t$ and $s$, except
  when $T = \theta(x)$ is an AC application of $f$ and $x$ occurs in a term $w$ that is not an AC
  application of $f$. Then $x$ is an argument both of $w$ and of $t$ or $s$, which have different
  keys, so $x\in\Vtwo(S_0)$, and the hypothesis says that $\theta(x)$ is a variable, a
  contradiction.
\end{proof}

\subsection{One AC round}

The step lemmas are composed along a round by induction on $\mathit{ac\_round}$.

\begin{lemma}[\isa{ac\_round\_meas}]
  Let $P$ consist of AC pairs, let all sides of $P \mathbin{@} R$ be non-variable terms and
  $\vars(P \mathbin{@} R) \subseteq \{0,\dots,\mathit{nv}-1\}$. For every result
  $(P_n, \sigma_n, \mathit{nv}_n)$ of $\mathit{ac\_round}\ P\ R\ \sigma\ \mathit{nv}$ there is a
  function $F$ on variables such that
  \begin{enumerate}[nosep]
    \item $\VNAC(P_n) \subseteq F(\VNAC(P \mathbin{@} R))$ and
          $\Vtwo(P_n) \subseteq F(\Vtwo(P \mathbin{@} R))$;
    \item $F(x) = x$ for every $x < \mathit{nv}$ not occurring in $P \mathbin{@} R$;
    \item every variable of $P_n$ occurs in $P \mathbin{@} R$ or is at least $\mathit{nv}$, and
          is below $\mathit{nv}_n$;
    \item if $\Vtwo(P_n) = F(\Vtwo(P \mathbin{@} R))$, there is $\rho$ with
          $\AS(P_n) \subseteq \AS(P \mathbin{@} R)\rho$ and
          $\sides(P_n) \subseteq (\AS(P \mathbin{@} R) \cup \sides(R))\rho$.
  \end{enumerate}
\end{lemma}

\begin{proof}[Proof sketch]
  $F$ is the composition of the renamings $\psi$ of the steps, and $\rho$ the composition of the
  substitutions $\theta$. The key point is item 4: we must derive the hypothesis of
  Lemma~\ref{lem:as-step} at each step from the equality at the end of the round. Suppose that at
  some step a variable $X \in \Vtwo(S_0)$ is bound to a non-variable term. Then $\psi(X) = X$, and
  $X$ does not occur in $S_1$, since $\theta$ is idempotent. By items 2 and 3 for the rest of the
  round, the later renamings leave $X$ unchanged and $X$ does not occur in $P_n$. Hence $X =
  F(X) \in F(\Vtwo(P \mathbin{@} R)) = \Vtwo(P_n) \subseteq \vars(P_n)$, a contradiction.
\end{proof}

\begin{theorem}[\isa{ac\_round\_mless}]
  If $P \neq []$ consists of AC pairs with variables below $\mathit{nv}$, every result $P_n$ of the
  round $\mathit{ac\_round}\ P\ []\ \sigma\ \mathit{nv}$ has a strictly smaller measure than $P$.
\end{theorem}

\begin{proof}
  By item 1, $|\VNAC(P_n)| \le |\VNAC(P)|$ and $|\Vtwo(P_n)| \le |\Vtwo(P)|$. If one of them is
  strict we are done. Otherwise $|\Vtwo(P_n)| = |\Vtwo(P)| \ge |F(\Vtwo(P))|$, and with item 1 this
  forces $\Vtwo(P_n) = F(\Vtwo(P))$; item 4 then gives $\rho$. If $P_n = []$, the last two
  components of its measure are $0$, while $\psize(P) > 0$. Otherwise, let $m$ be a side of $P_n$
  of maximal size. Then $m \in \AS(P)\rho$, but $m \notin \AS(P_n)$, since an admissible subterm
  of $P_n$ is strictly smaller than some side of $P_n$. Therefore
  $\AS(P_n) \subsetneq \AS(P)\rho$ and $|\AS(P_n)| < |\AS(P)\rho| \le |\AS(P)|$.
\end{proof}

\subsection{Termination of the algorithm}

\begin{theorem}[\isa{acunif\_dom}]
  For every problem $P$ with variables below $\mathit{nv}$ and every $\sigma$, the computation of
  $\mathit{acunif}\ P\ \sigma\ \mathit{nv}$ terminates. In particular
  $\mathit{ac\_unifiers}\ s\ t$ is defined for all terms $s$ and $t$.
\end{theorem}

\begin{proof}
  By well-founded induction on the measure, using the introduction rules of the domain predicate.
  Each syntactic step decreases the measure and preserves the bound on the variables; each result
  of an AC round has a smaller measure and its variables are below the new bound.
\end{proof}

%======================================================================
\section{Discussion}\label{sec:disc}
%======================================================================

\subsection{Main results}

For the Contejean--Devie solver, the final theorems read as follows, without any hypothesis.

\begin{corollary}
  Let $\mathit{ac\_unifiers}$ be instantiated with the Contejean--Devie solver. For all terms $s$
  and $t$:
  \begin{enumerate}[nosep]
    \item the computation of $\mathit{ac\_unifiers}\ s\ t$ terminates;
    \item every $\theta \in \mathit{ac\_unifiers}\ s\ t$ satisfies $s\theta \eqAC t\theta$;
    \item for every $\delta$ with $s\delta \eqAC t\delta$ there are
          $\theta \in \mathit{ac\_unifiers}\ s\ t$ and $\rho$ with $\delta(x) \eqAC \theta(x)\rho$
          for all $x \in \vars(s) \cup \vars(t)$.
  \end{enumerate}
\end{corollary}

Two immediate consequences, not stated separately in the formalization, are worth recording.
First, for a total solver, $s$ and $t$ are AC-unifiable if and only if
$\mathit{ac\_unifiers}\ s\ t \neq []$, which gives a decision procedure for AC-unifiability.
Second, AC-unification with any number of AC symbols and free symbols is finitary: the computed
list is a finite complete set of AC-unifiers. The computed set is not necessarily minimal.

\subsection{Main formal results}

\begin{center}
\small
\begin{tabular}{@{}p{0.48\textwidth}p{0.46\textwidth}@{}}
  \toprule
  Result & Isabelle name\\
  \midrule
  AC equality is decided by normal forms & \isa{eq\_AC\_iff\_nf}\\
  AC preserves sizes & \isa{eq\_AC\_size}\\
  Minimal solutions generate all solutions & \isa{Solutions\_gen\_by}\\
  Correctness of Contejean--Devie & \isa{cd\_solve\_correct}\\
  Termination of Contejean--Devie & \isa{cd\_terminates}\\
  Soundness of an AC step & \isa{solve\_ac\_sound}\\
  Soundness of the algorithm & \isa{acunif\_sound}\\
  Completeness of an AC step & \isa{solve\_ac\_complete}\\
  Completeness of the algorithm & \isa{acunif\_complete}\\
  Measure of one AC step & \isa{VNAC\_step}, \isa{V2\_step}, \isa{AS\_step}\\
  Measure of one AC round & \isa{ac\_round\_meas}, \isa{ac\_round\_mless}\\
  Termination of the algorithm & \isa{acunif\_dom}\\
  Final theorems & \isa{ac\_unifiers\_sound}, \isa{ac\_unifiers\_complete}\\
  Instance with Contejean--Devie & \isa{ac\_unifiers\_cd\_sound},\newline \isa{ac\_unifiers\_cd\_complete}\\
  \bottomrule
\end{tabular}
\end{center}

\subsection{Differences with the PVS development}

\begin{itemize}[leftmargin=*]
  \item \emph{Diophantine equations.} PVS enumerates candidate solutions up to a bound. Here the
        solver is an abstract parameter; soundness and termination hold for every solver, and
        completeness for every total one. The Contejean--Devie instance computes exactly the
        minimal solutions and is proved to terminate.
  \item \emph{Completeness.} Instead of renamings and restrictions of substitutions to the
        variables of the input, the proof keeps the invariant that an extension of the given
        unifier AC-fixes the computed substitution. This also removes the hypothesis, present in
        PVS, that the domain of the given unifier is contained in the variables of the problem.
  \item \emph{Signature.} Several AC symbols are given by a predicate. Since AC axioms only concern
        binary applications, other uses of AC symbols are free, and no well-formedness invariant
        on terms is needed.
  \item \emph{Axioms.} AC is defined by the set of all instances of its axioms, within a general
        inductive definition of equational theories, rather than by axioms over fixed variables.
  \item \emph{Partial definition.} The algorithm is defined as a partial function; soundness and
        completeness are proved by partial induction, and the termination theorem then
        discharges the domain conditions.
\end{itemize}

\subsection{Structure of the formalization}

\begin{center}
\begin{tabular}{@{}llr@{}}
  \toprule
  Theory & Content & Lines\\
  \midrule
  \isa{Equational\_Theory} & equational theories & 120\\
  \isa{AC\_Theory} & AC axioms of one symbol & 92\\
  \isa{AC\_Terms} & several AC symbols, normal forms, flattening & 468\\
  \isa{Dioph\_Solver} & interface to Diophantine solvers & 96\\
  \isa{Dioph\_CD} & Contejean--Devie algorithm and its correctness & 539\\
  \isa{Dioph\_CD\_Termination} & termination of Contejean--Devie & 548\\
  \isa{AC\_Unif\_Algorithm} & the unification algorithm & 206\\
  \isa{AC\_Unif\_Soundness} & soundness & 560\\
  \isa{AC\_Unif\_Completeness} & completeness & 951\\
  \isa{AC\_Unif\_Termination} & termination & 2\,596\\
  \isa{AC\_Unif\_CD} & instantiation with Contejean--Devie & 51\\
  \midrule
  Total & & 6\,227\\
  \bottomrule
\end{tabular}
\end{center}

\begin{thebibliography}{9}
\bibitem{AyalaRincon22}
M.~Ayala-Rinc\'on, M.~Fern\'andez, G.~F.~Silva, D.~Nantes-Sobrinho.
\newblock A certified algorithm for AC-unification.
\newblock In \emph{Formal Structures for Computation and Deduction (FSCD)}, 2022.

\bibitem{AyalaRincon24}
M.~Ayala-Rinc\'on, M.~Fern\'andez, G.~F.~Silva, T.~Kutsia, D.~Nantes-Sobrinho.
\newblock Certified first-order AC-unification and applications.
\newblock \emph{Journal of Automated Reasoning}, 2024.

\bibitem{BCD90}
A.~Boudet, E.~Contejean, H.~Devie.
\newblock A new AC unification algorithm with an algorithm for solving systems of Diophantine
equations.
\newblock In \emph{Logic in Computer Science (LICS)}, 1990.

\bibitem{CD94}
E.~Contejean, H.~Devie.
\newblock An efficient incremental algorithm for solving systems of linear Diophantine equations.
\newblock \emph{Information and Computation}, 113(1), 1994.

\bibitem{Fages87}
F.~Fages.
\newblock Associative-commutative unification.
\newblock \emph{Journal of Symbolic Computation}, 3(3), 1987.

\bibitem{Stickel81}
M.~E.~Stickel.
\newblock A unification algorithm for associative-commutative functions.
\newblock \emph{Journal of the ACM}, 28(3), 1981.

\bibitem{AFP-FOT}
C.~Sternagel, R.~Thiemann.
\newblock First-order terms.
\newblock \emph{Archive of Formal Proofs}, 2018.

\bibitem{AFP-HLDE}
F.~Me{\ss}ner, J.~Parsert, J.~Sch\"opf, C.~Sternagel.
\newblock Homogeneous linear Diophantine equations.
\newblock \emph{Archive of Formal Proofs}, 2017.

\bibitem{AFP-Farkas}
R.~Bottesch, M.~W.~Haslbeck, R.~Thiemann.
\newblock Farkas' lemma and Motzkin's transposition theorem.
\newblock \emph{Archive of Formal Proofs}, 2019.
\end{thebibliography}

\end{document}
